\section{Discussion}
\label{sec:discussion}
The matched reversals establish sensitivity beyond preserved summaries on
observed pairs; cross-corpus correlations cannot isolate this sensitivity
because size and topology also vary. Mostly small, localized and charge-dependent
responses, heterogeneous extraction, selected cycle frames and dependent edits
limit generalization. The measurements distinguish link arrangements without
identifying which serves readers better; that requires evidence about access
and use.

R0's first magnetic eigenvalue is lower than LT's
(Table~\ref{tab:refinement-summary}), despite LT's deliberate organization
around tasks and shared context; this does not establish worse LT information
flow. LT's full-graph density is slightly below the original's after adding pages
and links, whereas R0 is denser. The graphs omit semantic routes and link types;
modes can localize on few pages, and coverage and charge matter. Coupled text,
page and link edits, unequal area acceptance and one corpus per approach limit
interpretation. Metric-guided refinement requires accounting for these influences
and demonstrating useful access; a spectrum alone cannot identify useful links.

Agents used added links without clear primary answer separation; intervals
permit gains and losses. LT's lexical deficit shrinks under moved/derived-content
credit, showing that support attribution matters. Navigation does not establish
answer benefit.

Regression-test prioritization is a possible application, requiring comparison
with simpler change-impact and navigation measures where link traversal
supplies evidence beyond initial retrieval.
